\section{硬科技创业路线图：从技术男到产业链节点}

前面几章分别讲了降本、股权、融资、大单、定价和客户倒戈，本章把这些内容合成一张路线图。李一帆的故事里，技术能力只是第一步；公司必须逐步补上商业判断、客户理解、供应链能力、资本节奏和主机厂体系。硬科技创业者不是从实验室一步跳到产业巨头，而是在每一关里把一个短板补上。

\begin{knowledgebox}{硬科技创业阶段表}
\begin{tabularx}{\textwidth}{p{0.18\textwidth}p{0.34\textwidth}X}
\toprule
阶段 & 主要任务 & 典型风险 \\
\midrule
实验室阶段 & 证明技术可行 & 自嗨，找不到真实市场。 \\
样机阶段 & 让客户看见效果 & 只能演示，不能交付。 \\
订单阶段 & 获得真实预算 & 大单带来现金流和履约压力。 \\
量产阶段 & 降本、良率、供应链和售后 & 质量问题会摧毁信任。 \\
主机厂阶段 & 进入车型和长期供货体系 & 认证周期长，议价压力大。 \\
全球化阶段 & 面对不同资本和客户文化 & 信任语言和合规复杂。 \\
\bottomrule
\end{tabularx}
\end{knowledgebox}

\begin{importantbox}{硬科技创业者要补的课}
互联网创业者常先学增长和产品，硬科技创业者常先有技术。后者最难补的是客户、供应链、资本、量产和国际沟通。禾赛的 11 年史，就是一条不断补课的路线。
\end{importantbox}

\subsection{本章小结}

硬科技创业路线图说明，产业机会不是靠单点技术自动兑现。创业者要把技术推进到样机、订单、量产、主机厂和全球化，每一步都是新能力。

